\chapter*{Abstract}
\addcontentsline{toc}{chapter}{Abstract}

These notes provide a self-contained, primer-style summary of the
quantum-mechanical concepts most relevant for understanding modern
circuit quantum electrodynamics (cQED) experiments. We start from
the postulates of quantum mechanics, build up the formalism through
the two canonical model systems --- the quantum harmonic oscillator
and the two-level system --- and develop the tools needed to
describe open and driven systems: density matrices, time-dependent
perturbation theory, and the rotating-wave approximation. With this
foundation in place, we move to quantum optics, treating the
quantization of the electromagnetic field, Fock and coherent
states, and the Jaynes--Cummings model. The next parts specialise
to superconducting circuits, where we quantise lumped-element
circuits, introduce the Josephson junction, derive the standard
qubit Hamiltonians (transmon, fluxonium, \ldots), and discuss
readout, gates, decoherence, the basic calibration experiments,
and a short introduction to quantum error correction. The closing
part covers a research direction in which the same hardware
doubles as a single-quantum sensor: superconducting circuits as
detectors of feebly-coupled new physics, in particular sub-GeV
dark-matter recoils, dark photons, and axions.
